\documentclass[12pt]{article}
\usepackage[utf8]{inputenc}
\usepackage[T1]{fontenc}
\usepackage[margin=1in]{geometry}
\usepackage{setspace}
\usepackage{xcolor}
\usepackage[colorlinks=true, urlcolor=blue!60!black, linkcolor=blue!60!black]{hyperref}
\setlength{\parindent}{2em}
\setlength{\parskip}{1.25em}
\renewcommand{\baselinestretch}{1.0}
\begin{document}
\begin{center}
{\LARGE \textit{Dolan v. City of Tigard}, 512 U.S. 374 (1994)\par}
\vspace{0.8em}
{\large Case Brief --- Four-Page Edition\par}
\end{center}
\vspace{1em}

\section{Memory Jogger}

When Tigard conditioned Florence Dolan's store-expansion permit on dedicating a public floodplain greenway and a bike path, the Court held 5--4 that, beyond \textit{Nollan}'s essential nexus, the city must make an individualized determination of rough proportionality between the dedication and the development's impact, and it had not (\href{https://supreme.justia.com/cases/federal/us/512/374/}{Justia}).

\section{Detailed Case Facts}

Florence Dolan owned a 9,700-square-foot plumbing and electric supply store on a 1.67-acre parcel in downtown Tigard, Oregon. Fanno Creek crossed the parcel's corner, placing part of it in the 100-year floodplain. Tigard's code required 15\% open space in the business district, and the city had adopted a drainage plan preserving the floodplain as a greenway and a pedestrian/bicycle pathway plan. Dolan sought to nearly double the store to 17,600 square feet and pave a 39-space lot. The Planning Commission required her to dedicate (1) the floodplain portion as a public greenway and (2) an adjacent 15-foot strip for a public pathway---about 7,000 square feet, or 10\% of the parcel. It denied a variance, finding that more pavement ``can only add to the public need to manage the stream channel and floodplain'' and that a pathway ``could offset some of the traffic demand.''

\section{Procedural History}

Dolan argued that the dedications were an uncompensated taking. The Land Use Board of Appeals, the Oregon Court of Appeals, and the Oregon Supreme Court affirmed, finding a ``reasonable relationship'' between the conditions and the expansion's impacts. The Supreme Court granted certiorari.

\section{Judicial Votes}

Reversed and remanded, \textbf{5--4} (decided June 24, 1994).
\begin{itemize}
\item \textbf{Majority:} Chief Justice Rehnquist, joined by Justices O'Connor, Scalia, Kennedy, and Thomas.
\item \textbf{Dissents:} Justice Stevens, joined by Justices Blackmun and Ginsburg; Justice Souter.
\end{itemize}

\section{Holding}

Even when a permit condition satisfies \textit{Nollan}'s essential nexus, it violates the Takings Clause unless the government makes an individualized determination that the required dedication is roughly proportional---related in nature and extent---to the development's projected impact. Tigard did not show why flood control required a \emph{public} greenway rather than private open space, nor that the pathway would, rather than merely could, offset the store's added traffic.

\section{Analysis of Opinions}

\textbf{Majority (Rehnquist).} Demanding the strip outright would have been a per se taking of Dolan's right to exclude. Unlike citywide zoning upheld in \textit{Euclid}, Tigard made an adjudicative decision conditioning a permit on a deed of land, so the unconstitutional conditions doctrine applied. \emph{Step one:} flood control and traffic relief are legitimate purposes, and the conditions logically related to them, unlike the beach easement in \textit{Nollan}. \emph{Step two:} surveying state cases, the Court rejected both lax generalized-statement tests and Illinois's rigid specifically-and-uniquely-attributable test, choosing an intermediate standard and the term ``rough proportionality'' to avoid confusion with rational-basis review. The city bears the burden. \emph{Greenway:} the open-space rule already served flood control, and a permanent public easement stripped Dolan of all control over entry, unlike \textit{PruneYard}. \emph{Pathway:} despite an estimate of about 435 added daily trips, the city found only that a path could offset some demand.

\textbf{Stevens (dissenting).} State precedent did not support rough proportionality or shifting the burden to cities; the Court isolated the right to exclude instead of viewing the parcel as a whole and the benefits Dolan gained; ``could'' versus ``would'' was a ``play on words''; and heightened scrutiny of land-use permits revives \textit{Lochner}.

\textbf{Souter (dissenting).} The majority never really applied rough proportionality; both defects were \textit{Nollan} nexus problems, decided by improperly shifting the burden of proof.

\textbf{Later.} \textit{Lingle} (2005) preserved \textit{Nollan} and \textit{Dolan} as unconstitutional-conditions rules; \textit{Koontz} (2013) extended them to permit denials and money demands; and \textit{Sheetz} (2024) applied them to legislative exactions.

\section{Examples: Future Applications}

\begin{enumerate}
\item \textbf{Same-Side (Invalid): Boardwalk and Arts Fee.} A warehouse adding 15 truck trips a week must deed a public boardwalk easement and pay a \$250,000 arts fee; the fee lacks a nexus and the easement lacks proportionality.
\item \textbf{Same-Side (Invalid): Excessive Traffic Fee.} A \$60,000 fee on a backyard cottage that adds about \$3,000 of traffic burden is not roughly proportional.
\item \textbf{Opposite-Side (Valid): Private Buffer and Calibrated Fee.} A private stormwater buffer and a \$12,000 road fee based on an engineering study satisfy both tests.
\item \textbf{Opposite-Side (Valid): Turn Lane.} A traffic study shows a 300-home subdivision generates most peak traffic at its entrance, justifying a 12-foot turn-lane strip.
\item \textbf{Fence-Sitter: Citywide School Fee.} A per-square-foot school fee based on citywide averages, applied to a 55-and-older community that adds almost no students. \textit{Sheetz} left open whether class-based schedules must be tailored as specifically as parcel-by-parcel exactions; the city will cite its study, and the builder \textit{Dolan}'s demand for individualized findings.
\end{enumerate}

\section{Critique}

\textbf{Burden inversion.} Shifting the burden of proof to local governments and demanding quantified findings may lead cities to deny permits outright, where lenient \textit{Penn Central} review applies, or to rely on broad taxes rather than targeted mitigation fees.

\textbf{Tailoring legislative fees.} \textit{Sheetz} closed the legislative exemption but left acute tension with \textit{Dolan}'s demand for an individualized determination: how much aggregation may a citywide fee formula use?

\section{Key Quotations}

\begin{quote}
``We think a term such as `rough proportionality' best encapsulates what we hold to be the requirement of the Fifth Amendment. No precise mathematical calculation is required, but the city must make some sort of individualized determination that the required dedication is related both in nature and extent to the impact of the proposed development.''
\end{quote}
\noindent --- Chief Justice Rehnquist, for the Court, 512 U.S. at 391.

\begin{quote}
``Under the well-settled doctrine of `unconstitutional conditions,' the government may not require a person to give up a constitutional right---here the right to receive just compensation when property is taken for a public use---in exchange for a discretionary benefit conferred by the government where the benefit sought has little or no relationship to the property.''
\end{quote}
\noindent --- Chief Justice Rehnquist, for the Court, 512 U.S. at 385.

\begin{quote}
``We see no reason why the Takings Clause of the Fifth Amendment, as much a part of the Bill of Rights as the First Amendment or Fourth Amendment, should be relegated to the status of a poor relation in these comparable circumstances.''
\end{quote}
\noindent --- Chief Justice Rehnquist, for the Court, 512 U.S. at 392.

\end{document}
